\section{Implementierung des Patches}
\label{sec:implementierung}

In diesem Kapitel wird die Behebung des Sicherheitslecks vorgenommen.

\subsection{Anpassung von \texttt{get()} und \texttt{get\_or\_404()}}
\label{sec:patch-base}

Die zentrale Änderung betrifft ausschliesslich die Datei
\texttt{app/shared/services/base.py}.
Beide Methoden erhalten dieselben zwei optionalen Keyword-Only-Parameter, die
bereits in \texttt{get\_multi()} und \texttt{count()} vorhanden sind:

\begin{itemize}
    \item \texttt{organization\_id: str | None = None} \textendash{} wird
          nur ausgewertet, wenn das Modell das Attribut besitzt
          (\texttt{hasattr}-Guard, identisch zu \texttt{get\_multi()}).
    \item \texttt{include\_deleted: bool = False} \textendash{} filtert
          soft-gelöschte Datensätze, sofern das Modell \texttt{deleted\_at}
          kennt.
\end{itemize}

\texttt{get\_or\_404()} leitet beide Parameter an \texttt{get()} weiter und
liefert \texttt{404} zurück, wenn kein Treffer gefunden wird, also sowohl bei
nicht vorhandenen als auch bei organisationsfremden und soft-gelöschten
Datensätzen.

\texttt{delete()} wurde ebenfalls angepasst und ruft intern \texttt{get()} auf
und gibt \texttt{organization\_id} weiter, damit auch DELETE-Operationen auf
fremde Objekte ins Leere laufen.

Der vollständige Diff der drei Methoden ist nachfolgend dargestellt:

\begin{lstlisting}[style=custompython, caption={base.py \textendash{} gepatchte Methoden get(), get\_or\_404() und delete()}]
# VORHER
def get(self, db: Session, id: str) -> ModelType | None:
    result = db.execute(select(self.model).where(self.model.id == id))
    return result.scalar_one_or_none()

def get_or_404(self, db: Session, id: str) -> ModelType:
    obj = self.get(db, id)
    if obj is None:
        raise NotFoundError(self.model.__name__, id)
    return obj

# NACHHER
def get(self, db, id, *, organization_id=None, include_deleted=False):
    query = select(self.model).where(self.model.id == id)
    if organization_id and hasattr(self.model, "organization_id"):
        query = query.where(self.model.organization_id == organization_id)
    if not include_deleted and hasattr(self.model, "deleted_at"):
        query = query.where(self.model.deleted_at.is_(None))
    return db.execute(query).scalar_one_or_none()

def get_or_404(self, db, id, *, organization_id=None, include_deleted=False):
    obj = self.get(db, id,
                   organization_id=organization_id,
                   include_deleted=include_deleted)
    if obj is None:
        raise NotFoundError(self.model.__name__, id)
    return obj

def delete(self, db, id, soft_delete=True, *, organization_id=None):
    obj = self.get(db, id, organization_id=organization_id)
    if obj is None:
        return None
    # (Soft-/Hard-Delete-Logik unveraendert)
\end{lstlisting}

Die neuen Parameter sind keyword-only und optional, womit die
Signatur rückwärtskompatibel bleibt und das Refactoring der Aufrufer
pro Modul möglich wird. Der Mandantenfilter greift aber erst, sobald
ein Aufrufer \texttt{organization\_id} tatsächlich mitgibt. Die
Umstellung aller 78 Pflicht-Aufrufer zeigt der nächste Abschnitt.

\subsection{Anpassung aller Aufrufer}
\label{sec:patch-callers}

Alle 78 als Pflicht-Fix klassifizierten
\texttt{get\_or\_404()}-Aufrufer (Tabelle \ref{tab:callers_summary})
sowie die betroffenen \texttt{delete()}-Aufrufe werden nach einem
einheitlichen Muster angepasst.

\begin{lstlisting}[style=custompython, caption={Aufruf-Muster vor und nach dem Patch (Beispiel Fahrzeug-Endpunkt).}]
# VORHER -- kein Mandanten-Check
vehicle = vehicle_service.get_or_404(db, vehicle_id)

# NACHHER -- organisation_id wird mitgegeben
vehicle = vehicle_service.get_or_404(
    db, vehicle_id, organization_id=current_user.organization_id
)
\end{lstlisting}

Das gleiche Muster gilt für \texttt{delete()}-Aufrufe. Die fünf
Aufrufstellen in \texttt{shared/api/\allowbreak organizations.py}
bleiben unverändert, weil das \texttt{Organization}-Modell keinen
Mandantenschlüssel führt und die Endpunkte ausschliesslich für
Super-Admins zugänglich sind.

\paragraph{Service-interne Cross-Service-Lookups.}
Fünf Service-Methoden rufen intern \texttt{get\_or\_404()} einer
anderen Service-Instanz auf. Sie erhalten einen optionalen Parameter
\texttt{organization\_id}, der an das innere \texttt{get\_or\_404()}
durchgereicht wird: \texttt{dispense\_fuel()},
\texttt{receive\_\allowbreak delivery()}, \texttt{adjust\_level()},
\texttt{update\_maintenance()} der Klasse \texttt{FuelPumpService}
sowie \texttt{return\_tool()} der Klasse \texttt{ToolAssignmentService}.
Die aufrufenden API-Endpunkte befüllen den Parameter mit
\texttt{current\_user.\allowbreak organization\_id}.

\begin{lstlisting}[style=custompython, caption={Service-interner Cross-Service-Lookup mit Mandantencheck.}]
# fleet/services/fuel_pump_delivery.py -- create()
pump = fuel_pump_service.get_or_404(
    db, obj_in.pump_id, organization_id=obj_in.organization_id
)

# fleet/services/fuel.py -- create()
vehicle = vehicle_service.get_or_404(
    db, obj_in.vehicle_id, organization_id=obj_in.organization_id
)
\end{lstlisting}

\subsection{Behandlung von Modellen ohne \texttt{organization\_id}}
\label{sec:patch-no-org}

Der \texttt{hasattr}-Guard in \texttt{get()} stellt sicher, dass ein
Aufruf mit \texttt{organization\_id} bei Modellen ohne dieses Feld
stillschweigend ignoriert wird. Zwei dieser Modelle brauchen dennoch
eigene Regeln.

\begin{itemize}
    \item \texttt{Organization} ist der Mandant selbst. Die Endpunkte
          in \texttt{organizations.py} sind bereits ausschliesslich für
          Super-Admins zugänglich und bleiben unverändert.
    \item \texttt{Notification} ist nicht an eine Organisation, sondern
          über \texttt{user\_id} an den jeweiligen Benutzer gebunden.
          Der Schutz erfolgt deshalb über einen expliziten
          Eigentümercheck im Router (siehe
          Abschnitt~\ref{sec:patch-zusatzbefunde}, Absatz
          \emph{Notifications}).
\end{itemize}

\subsection{Fixes ausserhalb von \texttt{BaseService}}
\label{sec:patch-zusatzbefunde}

Die in Abschnitt \ref{sec:befunde-ausserhalb-baseservice} dokumentierten
Lücken betreffen Code, der \texttt{BaseService} umgeht oder nach einem
\texttt{get\_or\_404}-Lookup keinen Eigentümercheck vornimmt. Der
Wurzel-Fix greift dort nicht, weshalb diese Stellen einzeln behoben
werden. Alle Änderungen folgen dem bereits eingeführten Muster,
\texttt{current\_user.\allowbreak organization\_id} oder
\texttt{current\_user.\allowbreak id} wird dem Service oder einer
nachgelagerten Prüfung explizit übergeben.

\paragraph{Notifications \textendash{} Ownership-Check nach dem Lookup.}
Für Notifications genügt der Mandantenfilter nicht, weil sie an eine
\texttt{user\_id} gebunden sind. Nach dem \texttt{get\_or\_404}-Lookup
wird in GET, \texttt{POST /read} und DELETE geprüft, ob
\texttt{notification.user\_id == current\_user.id} gilt. Andernfalls
wird \texttt{404} zurückgegeben, um keine Information über fremde IDs
preiszugeben.

\paragraph{SAT-Assistances \textendash{} \texttt{get\_detail()} mit Mandanten-Parameter.}
\texttt{AssistanceService.get\_detail()} baut eigene SQLAlchemy-Filter.
Die Signatur erhält einen Parameter \texttt{organization\_id}, den der
Router mit \texttt{current\_user.\allowbreak organization\_id} befüllt.
Der zusätzliche \texttt{WHERE organization\_id = \$1} fliesst in alle
Teilabfragen der Methode ein.

\paragraph{Admin-Endpunkte in \texttt{users.py}.}
Die drei Router-Aufrufe von GET, PATCH und DELETE werden durch die im
vorherigen Abschnitt beschriebene Umstellung aller
\texttt{get\_or\_404}-/\texttt{delete}-Aufrufer mitgezogen, sobald
\texttt{organization\_id=\allowbreak current\_user.\allowbreak organization\_id}
ergänzt wird. Der zentrale Mandantenfilter in \texttt{BaseService.get()}
übernimmt damit auch hier die Prüfung.

\paragraph{DELETE in \texttt{sat/api/attachments.py}.}
\texttt{SatAttachmentService.delete()} nimmt bereits einen
\texttt{organization\_id}-Parameter entgegen, der Router übergab ihn
jedoch nicht. Der Router-Aufruf wird nach demselben Muster wie in
Abschnitt \ref{sec:patch-callers} um
\texttt{organization\_id=\allowbreak current\_user.\allowbreak organization\_id}
ergänzt, womit der im Service vorhandene Filter tatsächlich greift.

\paragraph{GPS \textendash{} Trip-Routen mit Mandantencheck.}
\begin{sloppypar}
Drei Trip-Routen wurden entlang desselben Musters abgesichert.
Bei \texttt{POST\allowbreak /trips/\allowbreak \{trip\_id\}/\allowbreak position}
wird vor dem Schreiben einer Position der zugehörige Trip über
\texttt{trip\_service.\allowbreak get\_or\_404(db,\allowbreak{} trip\_id,\allowbreak{}
organization\_id=current\_user.\allowbreak organization\_id)} geladen.
\texttt{POST\allowbreak /trips/\allowbreak \{trip\_id\}/\allowbreak end} und
\texttt{GET\allowbreak /trips/\allowbreak \{trip\_id\}/\allowbreak route}
laden den Trip ebenfalls über die gleiche \texttt{get\_or\_404}-Signatur,
bevor sie das Trip beenden bzw. die Routenpunkte zurückgeben. Gehört der
Trip nicht zum Mandanten des Aufrufers, brechen alle drei Endpunkte mit
\texttt{404} ab, bevor überhaupt gelesen oder geschrieben wird.
\end{sloppypar}

\paragraph{FK-List-Methoden \textendash{} Pflichtparameter
\texttt{organization\_id}.}
Alle 19 betroffenen List-Methoden erhalten einen
\textit{Pflicht}-Parameter \texttt{organization\_id: str} ohne Default.
Damit wird der Mandantenfilter zwingend an die Abfrage gehängt.
Vergisst ein Aufrufer den Parameter, reagiert Python mit
\texttt{TypeError} zur Laufzeit, was deutlich sicherer ist als ein
stilles Sicherheitsleck bei optionalem Parameter. Alle zugehörigen
Router befüllen den Parameter mit
\texttt{current\_user.\allowbreak organization\_id}.

\subsection{Zusammenfassung: Befund, Fix und Testabdeckung}
\label{sec:patch-uebersicht}

Tabelle \ref{tab:analyse_patch_test} verknüpft jeden Befund aus der
Analyse mit dem zugehörigen Patch-Unterabschnitt und dem Test, der den
Fix absichert. Die Smoke- und Integrationstests der bestehenden Suite
dienen in allen Fällen als zusätzlicher Regressionsschutz.

\begin{table}[H]
      \begin{center}
            \renewcommand{\arraystretch}{1.3}
            \begin{tabular}{|p{4.8cm}|l|l|p{3.8cm}|}
                  \hline
                  \textbf{Befund (Issue)}
                   & \textbf{Analyse}
                   & \textbf{Fix}
                   & \textbf{Test}                                                                                   \\
                  \hline
                  IDOR in \texttt{get()} / \texttt{get\_or\_404()} (\#5)
                   & \ref{sec:analyse}
                   & \ref{sec:patch-base}
                   & Angriffsklasse 1                                                                                \\
                  Soft-Delete-Bypass in \texttt{get()} / \texttt{get\_or\_404()} (\#14)
                   & \ref{sec:analyse}
                   & \ref{sec:patch-base}
                   & Angriffsklasse 2                                                                                \\
                  78 Router-Aufrufer ohne \texttt{organization\_id}
                   & \ref{sec:analyse}
                   & \ref{sec:patch-callers}
                   & \texttt{TestIDOR}, \texttt{TestIDOREmployee}                                                    \\
                  5 Service-interne Cross-Service-Aufrufe
                   & \ref{sec:analyse}
                   & \ref{sec:patch-callers}
                   & Smoke-Tests \texttt{fuel}/\texttt{tool}                                                         \\
                  Notifications ohne Ownership-Check
                   & \ref{sec:befunde-ausserhalb-baseservice}
                   & \ref{sec:patch-zusatzbefunde}
                   & Smoke-Tests \texttt{notifications}                                                              \\
                  SAT-Assistances ohne Mandantenfilter
                   & \ref{sec:befunde-ausserhalb-baseservice}
                   & \ref{sec:patch-zusatzbefunde}
                   & Smoke-Tests \texttt{sat}                                                                        \\
                  GPS-Position ohne Trip-Ownership-Check
                   & \ref{sec:befunde-ausserhalb-baseservice}
                   & \ref{sec:patch-zusatzbefunde}
                   & Smoke-Tests \texttt{gps}                                                                        \\
                  19 FK-List-Methoden ohne Mandantenfilter
                   & \ref{sec:befunde-ausserhalb-baseservice}
                   & \ref{sec:patch-zusatzbefunde}
                   & Smoke-/Integrationstests                                                                        \\
                  \hline
            \end{tabular}
      \end{center}
      \caption{Zuordnung jedes Befunds zu seinem Patch-Abschnitt und der Testabdeckung.}
      \label{tab:analyse_patch_test}
\end{table}
